La primera respuesta generada por la IA produjo un bloque \texttt{flowchart LR} sintácticamente válido, precedido de un texto introductorio, que renderizó sin errores en mermaid.live y respetaba las reglas de la plantilla: una única asociación actor--caso de uso, estereotipos con comillas latinas y \texttt{classDef}/\texttt{class} para el actor. No obstante, tras auditar el código frente a la semántica UML, a los criterios de aceptación de US-6 y al glosario común, se aplicaron las siguientes correcciones y mejoras:
\begin{itemize}[leftmargin=*,nosep]
\item \textbf{Uso incorrecto de «include» para filtros opcionales:} La IA modeló «Filtrar por categoría» y «Filtrar por talla» como relaciones «include» de «Filtrar catálogo». El cliente puede filtrar solo por una de ellas o por ninguna, por lo que no son pasos obligatorios. Se convirtieron en relaciones «extend» discontinuas dirigidas desde el caso opcional hacia el caso base, y se añadió «Limpiar filtros» como extensión adicional (coherente con el flujo opcional del diagrama de secuencia de US-6).
\item \textbf{Relación «extend» invertida respecto a la semántica del criterio:} La IA modeló «Mostrar contador de resultados» como extensión de «Filtrar dinámicamente sin recargar», pero el criterio de aceptación exige que el contador se muestre \emph{siempre}. Se convirtió en un caso incluido por «Filtrar catálogo», junto con «Actualizar listado sin recargar la página».
\item \textbf{Requisito no funcional modelado como caso de uso:} «Filtrar dinámicamente sin recargar» describe una cualidad del sistema (actualización sin recarga completa), no un objetivo del actor. Se reformuló como el caso de uso incluido «Actualizar listado sin recargar la página».
\item \textbf{Nombre de actor y estilo inconsistentes con el glosario:} La IA usó el actor \texttt{Clienta} y nombres con mayúsculas iniciales en cada palabra, distintos del glosario común (\texttt{Cliente}) y del resto del proyecto. Se unificó el actor y el estilo de los nombres, se usaron identificadores descriptivos en lugar de siglas (\texttt{FC}, \texttt{FCA}, \texttt{FT}, \texttt{FD}, \texttt{MC}) y el límite pasó a titularse «Atelier Jaén - Filtros del catálogo».
\end{itemize}
